\documentclass[a4paper, 12pt]{article}
\usepackage{pdfpages}
\usepackage[T2A]{fontenc}
\usepackage[utf8]{inputenc}
\usepackage[russian]{babel}
\usepackage{mathtext}
\usepackage{amsmath, amssymb}
\usepackage{graphicx}
\usepackage{listings}
\usepackage{hyperref}
\usepackage{revsymb}
\usepackage{minted}
\usepackage{tabularx} % Для равномерной ширины столбцов 
\usepackage{geometry} % Для настройки полей
\usepackage{amsthm} % Для создания окружений типа "определение"
\newtheorem{definition}{Определение}

\hypersetup{
    colorlinks=true,
    linkcolor=blue,
    filecolor=magenta,
    urlcolor=cyan,
}

\urlstyle{same}
\DeclareGraphicsExtensions{.pdf,.png,.jpg,.jpeg}
\setcounter{MaxMatrixCols}{20}
\graphicspath{{pictures/}}

%\geometry{
%    a4paper,
%    left=20mm,   % Левое поле
%    right=20mm,  % Правое поле
%    top=20mm,    % Верхнее поле
%    bottom=20mm  % Нижнее поле
%}


\title{\textbf{Комментарии к основаниям математики \\
	Commentary on the Foundations of Mathematics}}
\author{И.А.Юхновский}
\date{Сентябрь 2025}
    

\begin{document}
\maketitle
\thispagestyle{empty}
\section*{Аннотация}
Комментарии к основаниям математики

\section*{Abstract}
Commentary on the Foundations of Mathematics

\section{Субъект (Subject)}
Субъект — это то, о чём что-то утверждается в предложении.
В классической логике субъект — это носитель свойств или отношений.

В контексте Principia Mathematica субъект часто ассоциируется с переменными или индивидуальными объектами, которые могут быть аргументами функций или предикатов.

Например, в высказывании \textit{Socrates — человек}  \textit{Socrates — субъект}.

\section{Объект (Object)}
Объект — это то, что утверждается о субъекте, или то, на что направлено действие/отношение.
В Principia Mathematica объект может быть:
\begin{itemize}
\item  Конкретным индивидуальным объектом (например, число, множество).
\item  Функцией, предикатом или отношением (например, свойство \textit{быть человеком}).
\end{itemize}

\section{Ключевое различие между субъектом и объектом}
\begin{itemize} 
\item  Субъект — это то, что подставляется в логическое высказывание (например, переменная или индивид).
\item  Объект — это то, о чём или с чем оперирует высказывание (например, свойство, множество, функция).
\end{itemize} 

В теории типов Рассела и Уайтхеда объекты делятся на иерархию типов, чтобы избежать самоприменимости (например, "множество всех множеств, не содержащих себя").

\section{Примеры}

\subsection{Простое высказывание (предикатная логика)}
\subsubsection{Пример 1}
Высказывание: x принадлежит множеству A 

В математической записи:
\[x \in A\]

В формальной записи:
\begin{minted}{haskell}
A(x)
\end{minted}

Здесь:
\begin{itemize} 
 \item \textit{x} - субъект (индивидуальный объект). 
 \item \textit{A} - объект (множество, которое является объектом, о котором утверждается принадлежность).
\end{itemize} 

\subsubsection{Пример 2}
Высказывание:\textit{Сократ — человек}

В формальной записи:
\begin{minted}{haskell}
Человек(Сократ)
\end{minted}

Здесь:
\begin{itemize} 
 \item \textit{Сократ} — субъект (индивидуальный объект, о котором идёт речь).
 \item \textit{Человек} — объект (предикат, который приписывается субъекту).
\end{itemize} 

\subsection{Отношения между объектами}
Высказывание: \textit{Афины больше Спарты}

В формальной записи:
\begin{minted}{haskell}
Больше(Афины, Спарта)
\end{minted}

Здесь:
\begin{itemize} 
 \item \textit{Афины} и \textit{Спарта} — субъекты (аргументы).
 \item \textit{Больше} — объект
\end{itemize} 

\section{Почему это важно?}
Разделение на субъект и объект помогает:
\begin{itemize} 
	\item Избежать парадоксов (например, парадокса Рассела: "множество всех множеств, не содержащих себя").

	\item Построить строгую логическую систему, где каждый объект имеет определённый тип и не может быть сам себе субъектом и объектом одновременно.
\end{itemize} 

\section{Парадокс Рассела}
Формально парадокс Рассела можно записать следующим образом:
Обозначим через $R$ множество всех множеств, которые не содержат сами себя в качестве элемента:
\[R = \lbrace x | x \notin x \rbrace\]
Теперь зададим вопрос: принадлежит ли $R$ само себе?

\begin{itemize} 
 \item Если $R \in R$, то по определению $R$ должно удовлетворять условию $R \notin R$.
 \item Если $R \notin R$, то $R$ удовлетворяет условию $R \in R$.
\end{itemize} 

Таким образом, оба утверждения приводят к противоречию:

\[ R \in R \Leftrightarrow R \notin R \]

\subsection{Решение Рассела и Уайтхеда}
Ввести иерархию типов, чтобы субъект и объект всегда принадлежали разным типам. 

Например:
\begin{itemize} 
 \item Множество $x$ — тип 0.
 \item Множество множеств (например, $R$) — тип 1.
 \item Множество множеств множеств — тип 2, и так далее.
\end{itemize} 

Теперь $x \in x$ запрещено, так как $x$ не может быть одновременно субъектом и объектом одного и того же типа.

\section{Атомарные предложения}
Список вариантов из необработанных предложений, требуемых в доказательствах, может быть сведен к списку истинных атомарных предложений вместе с положением о том, что каждое истинное атомарное предложение есть одно из следующих:
\begin{itemize} 
 \item \textit{p},
 \item \textit{q},
 \item \textit{r},
 \item \textit{s},
 \item \textit{t}.
\end{itemize} 

\subsection{Типы атомарных предложений}
\begin{itemize} 
 \item $R_1(x)$ - субъект $x$ имеет предикатом $R_1$
 \item $R_2(x,y)$ - субъект $x$ находится в отношении $R_2$ (интенсионально) с субъектом $y$
 \item $R_3(x,y,z)$ - субъекты $x,y,z$ находятся в трехместном тернарном отношении $R_3$ (интенсионально)
  \item $R_4(x,y,z,w)$ - субъекты $x,y,z,w$ находятся в четырехместном тернарном отношении $R_4$ (интенсионально)
\end{itemize}

\subsection{Интенсионально}
В логике и философии термин «интенсионально» (от англ. intensional) относится к контекстам, в которых значение выражения зависит не только от его экстенсионала (объекта, на который оно указывает), но и от способа задания этого объекта, его смысла или описания.

\begin{table}[h]
    \centering
    \caption{Интенсиональность vs. Экстенсиональность}
    \label{tab:equalwidth}
    \begin{tabularx}{\textwidth}{|X|X|}
        \hline
        \textbf{Интенсиональность} & \textbf{Экстенсиональность} \\
        \hline
        Значение зависит от смысла, описания или способа представления объекта. & Значение зависит только от объекта, на который ссылается выражение. \\
        \hline
        Примеры: убеждения, желания, модальные операторы ("возможно", "необходимо"), временные контексты ("вчера", "завтра"). & Примеры: простые утверждения о фактах ("Москва — столица России"). \\
        \hline
    \end{tabularx}
\end{table}


\subsubsection{Примеры интенсиональных контекстов}


\textbf{Убеждения и знания:}

\begin{itemize} 
 \item \textit{Он верит, что утренняя звезда — это Венера.}
 \item \textit{Он верит, что вечерняя звезда — это Венера.}
\end{itemize} 

В обоих случаях речь идёт об одном и том же объекте (Венере), но убеждения могут различаться, если \textit{Он} не знает, что \textit{утренняя звезда} и \textit{вечерняя звезда} — это одно и то же.

\textbf{Модальные операторы:}

\begin{itemize} 
 \item \textit{Необходимо, что 2 + 2 = 4.} (Истинно)
 \item \textit{Необходимо, что количество планет в Солнечной системе — 8.} (Ложно, так как это факт, а не логическая необходимость)
\end{itemize} 


\textbf{Контрафактуальные утверждения:}

\begin{itemize} 
 \item \textit{Если бы я был президентом, я бы изменил законы.} (Истина зависит от контекста, а не от реального положения дел)
\end{itemize} 


\section{Молекулярные предложения}
\textit{p,q,r,s,t} - простые пропозициональные буквы.
\subsection{Пропозиция}
Пропозиция — это логическое суждение, которое может быть истинным или ложным.

Пропозиции являются основными "кирпичиками" логики и математики.
Они не зависят от языка, на котором выражены (в отличие от предложений естественного языка).

В современной логике пропозиции соответствуют высказываниям в логике высказываний или формулам в логике предикатов. Они лежат в основе формальных систем, используемых в математике, информатике и философии.

\begin{itemize} 
 \item \textit{2 + 2 = 4} — это пропозиция (истинная).
 \item \textit{Снег зелёный} — это пропозиция (ложная).
 \item \textit{Закрой дверь!} — не пропозиция, так как это не суждение, а команда.
\end{itemize}

\subsection{p|q}
$\mid$ - штрих-символ Шеффера.
\begin{table}[h]
    \centering
    \caption{Таблица истинности для $p \mid q$}
    \label{tab:equalwidth}
    \begin{tabularx}{\textwidth}{|X|X|X|}
        \hline
        \textbf{$p$} & \textbf{$q$} & \textbf{$p \mid q$}\\
        \hline
        0 & 0 & 1 \\
        \hline
        0 & 1 & 1 \\
        \hline
        1 & 0 & 1 \\
        \hline
        1 & 1 & 0 \\
        \hline
    \end{tabularx}
\end{table}

Истинно, когда хотя бы одно (или оба одновременно) из атомарных высказываний будет ложным.

$p \mid q$ означает, что:
\begin{itemize} 
 \item \textit{p несовместно q}
 \item \textit{p ложно или q ложно}
 \item \textit{p влечет не q} 
 \item \textit{p штрих q} 
\end{itemize}

%\begin{equation}
%
%\end{equation}

\subsection{Отрицание $\sim p$}
\begin{definition}[Df]
    \[
       \sim p.=.p \mid p 
    \]
\end{definition}

\begin{itemize} 
 \item \textit{Не p}
 \item \textit{Отрицание p равно по определению штриху Шеффера p с самим собой.}
\end{itemize}

\begin{table}[h]
    \centering
    \caption{Таблица истинности для $p \sim p $}
    \label{tab:equalwidth}
    \begin{tabularx}{\textwidth}{|X|X|X|}
        \hline
        \textbf{$p$} & \textbf{$\sim p$} & \textbf{$p \mid p$}\\
        \hline
        0 & 1 & 1 \\
        \hline
        1 & 0 & 0 \\
        \hline
    \end{tabularx}
\end{table}

\subsection{Импликация $p \supset q$}

\begin{definition}[Df]
    \[
       p \supset q .=. p \mid \sim q 
    \]
\end{definition}

\begin{itemize} 
 \item \textit{Если p, то q}
 \item \textit{p влечёт q}
\end{itemize}
 
Если утверждение $p$ истинно, то утверждение $q$ также должно быть истинным.

В классической логике импликация ложна только в одном случае: когда $p$ истинно, а $q$ ложно.

\begin{table}[h]
    \centering
    \caption{Таблица истинности для $p \supset q$}
    \label{tab:equalwidth}
    \begin{tabularx}{\textwidth}{|X|X|X|}
        \hline
        \textbf{$p$} & \textbf{$q$} & \textbf{$p \supset q$}\\
        \hline
        0 & 0 & 1 \\
        \hline
        0 & 1 & 1 \\
        \hline
        1 & 0 & 0 \\
        \hline
        1 & 1 & 1 \\
        \hline
    \end{tabularx}
\end{table}

\begin{itemize} 
 \item 1 — это "истина" (англ. true). Это означает, что высказывание истинно, верно или выполняется.
 \item 0 — это "ложь" (англ. false). Это означает, что высказывание ложно, неверно или не выполняется.
\end{itemize}

\subsubsection{Примеры}

\begin{itemize} 
 \item Если $p$ — \textit{идёт дождь}, а $q$ — \textit{земля мокрая}, то $p \supset q$ читается как \textit{Если идёт дождь, то земля мокрая.}

 \item В математических доказательствах $p \supset q$ может означать \textit{Из p следует q}.
 
\end{itemize}

\subsection{Дизъюнкция $p \vee q$}

\begin{definition}[Df]
    \[
       p \vee q .=. \sim p \mid \sim q 
    \]
\end{definition}

$p \vee q$ — это дизъюнкция (логическое \textit{ИЛИ}.

Читается как \textit{p или q}.

\begin{table}[h]
    \centering
    \caption{Таблица истинности для $p \vee q$}
    \label{tab:equalwidth}
    \begin{tabularx}{\textwidth}{|X|X|X|}
        \hline
        \textbf{$p$} & \textbf{$q$} & \textbf{$p \vee q$}\\
        \hline
        0 & 0 & 0 \\
        \hline
        0 & 1 & 1 \\
        \hline
        1 & 0 & 1 \\
        \hline
        1 & 1 & 1 \\
        \hline
    \end{tabularx}
\end{table}

Истинно, если хотя бы одно из высказываний \textit{p} или \textit{q} истинно.

\subsection{Конъюнкция $p.q$}

\begin{definition}[Df]
    \[
       p.q .=. \sim (p \mid q) 
    \]
\end{definition}

$p.q$ — это конъюнкция (логическое \textit{И}).

Читается как "p и q".

\begin{table}[h]
    \centering
    \caption{Таблица истинности для $p.q$}
    \label{tab:equalwidth}
    \begin{tabularx}{\textwidth}{|X|X|X|}
        \hline
        \textbf{$p$} & \textbf{$q$} & \textbf{$p.q$}\\
        \hline
        0 & 0 & 0 \\
        \hline
        0 & 1 & 0 \\
        \hline
        1 & 0 & 0 \\
        \hline
        1 & 1 & 1 \\
        \hline
    \end{tabularx}
\end{table}

Истинно, только если оба высказывания $p$ и $q$ истинны.

\subsection{Импликация $p \supset q$}

\begin{definition}[Df]
    \[
       p \supset q .=. p \mid (q \mid q) 
    \]
\end{definition}

\section{Примитивное предложение 1}
На основании $p$ и $p \mid ( g \mid r)$ вывести $r$

\subsection*{Определение штриха Шеффера}
Штрих Шеффера \( A | B \) определяется как:
\[
A | B = \neg (A \land B)
\]

\subsection*{Дано}
\[
p \quad \text{и} \quad p | (g | r)
\]

\subsection*{Цель}
Найти \( r \).

\subsection*{Логический анализ}

1. Раскроем выражение \( p | (g | r) \):
\[
p | (g | r) = \neg (p \land (g | r))
\]

2. Подставим определение штриха Шеффера для \( g | r \):
\[
g | r = \neg (g \land r)
\]

3. Подставим \( g | r \) в исходное выражение:
\[
p | (g | r) = \neg (p \land \neg (g \land r))
\]

4. Упростим выражение:
\[
p | (g | r) = \neg p \lor (g \land r)
\]

\subsection*{Анализ случаев}

\subsubsection*{Случай 1: \( p = 1 \) (истина)}
\[
\neg p = 0
\]
Тогда выражение упрощается до:
\[
0 \lor (g \land r) = g \land r
\]
Для истинности \( p | (g | r) \) необходимо, чтобы \( g \land r \) было истинно.

\subsubsection*{Случай 2: \( p = 0 \) (ложь)}
\[
\neg p = 1
\]
Тогда выражение упрощается до:
\[
1 \lor (g \land r) = 1
\]
В этом случае \( p | (g | r) \) всегда истинно, независимо от \( g \) и \( r \).

\subsection*{Вывод}
Если \( p \) истинно (\( p = 1 \)), то для истинности \( p | (g | r) \) необходимо, чтобы \( g \land r \) было истинно. Это означает, что \( r \) должно быть истинно, если \( g \) истинно.

\section{Примитивное предложение 1}
$\lbrace (p \mid q) \mid r \rbrace \mid s .\equiv :. \sim p \vee q.r: \vee : \sim s$

$a \mid b$ — это штрих Шеффера, который эквивалентен $\sim (a \wedge b)$.

\subsection*{Шаг 1: Преобразуем внутренние операции}
\[(p \mid q) \equiv \sim (p \wedge q)\]

\subsection*{Шаг 2: Подставляем результат в следующую операцию}

\[\lbrace (p \mid q) \mid r \rbrace \equiv \sim \left( \sim (p \wedge q) \wedge r \right) \]

Используем закон де Моргана:
\[\sim \left( \sim (p \wedge q) \wedge r \right) \equiv \sim (\sim (p \wedge q)) \vee \sim r \equiv (p \wedge q) \vee \sim r\]

\subsection*{Шаг 3: Подставляем результат в последнюю операцию}
\[\lbrace (p \mid q) \mid r \rbrace \mid s \equiv \sim \left( (p \wedge q) \vee \sim r \wedge s \right)\]

Используем закон де Моргана ещё раз:
\[\sim \left( (p \wedge q) \vee \sim r \wedge s \right) \equiv \sim (p \wedge q) \wedge \sim (\sim r \wedge s)\]

\subsection*{Шаг 4: Упрощаем выражение}
\[\sim (p \wedge q) \equiv \sim p \vee \sim q\]
\[\sim (\sim r \wedge s) \equiv r \vee \sim s\]

Объединяем результаты:
\[(\sim p \vee \sim q) \wedge (r \vee \sim s)\]

\subsection*{Шаг 5: Раскрываем скобки}
\[(\sim p \vee \sim q) \wedge (r \vee \sim s) \equiv (\sim p \wedge r) \vee (\sim p \wedge \sim s) \vee (\sim q \wedge r) \vee (\sim q \wedge \sim s)\]

\subsection*{Шаг 6: Сравниваем с правой частью}
Правая часть выражения:
\[ \sim p \vee q \cdot r \vee \sim s\]

Чтобы показать эквивалентность, заметим, что:

$\sim p \vee \sim s$ присутствует в обеих частях.
$q \cdot r$ соответствует $\sim q \wedge r$ через отрицание.

Таким образом, после упрощения обе части эквивалентны.
\subsection*{Итого:}
\[ \lbrace (p \mid q) \mid r \rbrace \mid s \equiv \sim \left( \sim \left( \sim (p \wedge q) \wedge r \right) \wedge s \right) \]
\[\equiv \sim \left( (p \wedge q) \vee \sim r \right) \wedge s \]
\[\equiv (\sim p \vee \sim q) \wedge (r \vee \sim s) \]
\[\equiv (\sim p \wedge r) \vee (\sim p \wedge \sim s) \vee (\sim q \wedge r) \vee (\sim q \wedge \sim s) \]
\[\equiv \sim p \vee q \cdot r \vee \sim s\]

\begin{thebibliography}{9}
    \bibitem{russell1910principia}
    Russell, B., \& Whitehead, A. N. (1910--1913).
    \textit{Principia Mathematica}.
    Cambridge University Press.
\end{thebibliography}

\end{document}